% T2.8 / R1-4 — "Have the authors considered FP8, FP4 and such number formats
% as they seem to be more popular than traditional fixed point implementations
% especially in LLMs."
%
% Place at the end of Section III, after the synthesis methodology.
%
% The answer is not a deflection: the precision sweep we already ran quantifies
% exactly what an FP8-class mantissa would do to the rotation, and the
% block-normalisation that S-14 forced on us is already the floating-point-like
% scheme that suits this operation. All numbers below are measured.

\subsection{On FP8, FP4 and Block Floating Point}
\label{sec:formats}

Low-precision floating-point formats have become the default for LLM
inference, so it is reasonable to ask why a rotation datapath should be fixed
point. Three observations, two of them quantitative.

\subsubsection{The formats target a different bottleneck}
FP8 and FP4 are deployed against weight storage and the general matrix
multiplies that dominate transformer inference --- operations whose cost is
memory traffic and multiplier count, and whose operands span a wide dynamic
range. RoPE is neither. It is a fixed, data-independent rotation applied to
activations already resident in registers, and it contains no weights at all.
Narrowing the format of a rotation does not relieve the bottleneck those
formats were designed for.

\subsubsection{The datapath here is already floating point, in the form that fits}
As Section~\ref{sec:frontend} describes, a bare fixed-point rotator is not
deployable: measured query and key activations reach $|q| = 132$ and
$|k| = 178$, far outside the $[-2,2)$ range of the $Q(1,F)$ format, so the
inputs are block-normalised with a shared exponent per (head, token) before
the rotator and de-normalised after. That is block floating point, and it is
the appropriate choice for this operation specifically. The $d_h/2$ pairs of
one head at one position share a magnitude scale closely enough that a single
exponent serves them all. We measured what a finer exponent is actually worth
by sweeping the block size over real query activations captured from a Qwen2
model (Table~\ref{tab:blocksize}). Moving from one exponent per head-token to
a per-element exponent --- which is what an FP8 or FP4 element format provides
--- reduces the rotation error by roughly $3.8\times$, that is, it buys about
two bits of mantissa. Beyond a block of $16$ pairs there is no further gain at
all: the exponent is a power of two, and block maxima rarely fall in different
octaves.

\begin{table}[t]
\caption{Relative rotation error against block size $D$ (pairs per shared
exponent), Binary UD $N=10$, on real Qwen2 query activations. $D=64$ is one
exponent per head-token; $D=1$ is a per-element float.}
\label{tab:blocksize}
\centering
\resizebox{\columnwidth}{!}{%
\begin{tabular}{crrrrrr}
\toprule
$F$ & $D{=}1$ & $D{=}2$ & $D{=}4$ & $D{=}8$ & $D{=}16$ & $D{=}64$ \\
\midrule
14 & 0.00034 & 0.00036 & 0.00038 & 0.00042 & 0.00048 & 0.00057 \\
12 & 0.00053 & 0.00067 & 0.00090 & 0.00120 & 0.00158 & 0.00204 \\
11 & 0.00088 & 0.00126 & 0.00186 & 0.00258 & 0.00341 & 0.00445 \\
10 & 0.00178 & 0.00279 & 0.00397 & 0.00558 & 0.00750 & 0.00993 \\
\bottomrule
\end{tabular}%
}
\end{table}

Two bits is not free, but neither is a per-element exponent: it adds an
exponent field, alignment and normalisation to every lane in place of one
reduction tree and leading-zero count per block. This design takes the coarse
block, paying two bits of datapath width to save that logic on $63$ of every
$64$ lanes.

\subsubsection{The precision sweep already answers the question}
What the rotation needs is angular resolution on a normalised operand, which
is mantissa precision, not exponent range. Section~\ref{sec:precision}
measures that requirement directly: $F = 11$ fractional bits for perplexity
indistinguishable from floating point across five models, and a collapse below
$F = 9$ for a model with a wide frequency schedule. An E4M3 FP8 value carries
three explicit mantissa bits, and FP4 carries one. Placed on the same axis,
those sit well past the cliff --- and our sweep records what happens there:
at $F = 4$, LLaMA-2-7B degrades by $+4.9\times10^{4}\%$ and Qwen2-7B by
$+2.4\times10^{5}\%$.

The reason is the mechanism of \S\ref{sec:precision}: once the rotation angle
for a dimension pair falls below the datapath LSB, that pair ceases to rotate
and its positional information is discarded rather than approximated. A format
that spends its bits on exponent range cannot help, because the operand has
already been normalised and the range was never the problem.

The conclusion is not that low-precision floating point is unsuited to LLM
inference --- it plainly is suited --- but that the rotation is the wrong
place to apply it. A shared-exponent fixed-point datapath at $F = 11$ costs
less silicon than an FP8 lane of equivalent mantissa precision would, and the
measurements in Table~\ref{tab:perf} are taken at that operating point.
